\chapter{Propositional Logic}
\signature

\section{Propositions}

\begin{definition}[Proposition]
A \emph{proposition} is a declarative sentence that is either true ($\TT$) or
false ($\FF$), but never both and never neither. Propositions are denoted by
lowercase letters $p, q, r, \ldots$ and take their values in the set
$\B = \{\TT, \FF\}$ of \emph{truth values}.
\end{definition}

\begin{example}
``$2+2=4$'' is a true proposition; ``Paris is the capital of Spain'' is a false
proposition. By contrast, ``Close the door'' (imperative), ``What time is it?''
(interrogative) and ``$x + 1 = 0$'' (free variable --- a \emph{predicate},
Chapter~4) are not propositions. The self-referential sentence ``This sentence
is false'' is also excluded: it cannot consistently receive either truth value,
violating the principle of bivalence.
\end{example}

\begin{definition}[Atomic and compound]
An \emph{atomic} proposition cannot be decomposed using logical connectives; it
is symbolised by a single letter. A \emph{compound} proposition is built from
atomic propositions using connectives such as
$\neg, \wedge, \vee, \to, \leftrightarrow$.
\end{definition}

Propositional logic studies only how the truth of a compound proposition
depends on its atoms; looking \emph{inside} an atom is the job of predicate
logic (Chapter~4).

\subsection{Well-formed formulas}

The propositional language can be defined by a grammar: every atomic letter is
a well-formed formula (wff); if $\varphi$ and $\psi$ are wffs, so are
$\neg\varphi$, $(\varphi \wedge \psi)$, $(\varphi \vee \psi)$,
$(\varphi \to \psi)$ and $(\varphi \leftrightarrow \psi)$; nothing else is a
wff. This inductive definition (a recursive set in the sense of Chapter~14)
makes ``formula'' a precise mathematical object.

\begin{definition}[Truth assignment]
A \emph{truth assignment} (or \emph{valuation}) is a function
$v : \{p, q, r, \ldots\} \to \B$ assigning a truth value to every atomic
proposition. A valuation extends uniquely to all wffs via the truth tables of
the connectives.
\end{definition}

\section{The logical connectives}

\begin{definition}[Negation]
The \emph{negation} $\neg p$ (``not $p$'') is true exactly when $p$ is false.
\end{definition}

\begin{definition}[Conjunction and disjunction]
The \emph{conjunction} $p \wedge q$ (``$p$ and $q$'') is true exactly when both
$p$ and $q$ are true. The \emph{disjunction} $p \vee q$ (``$p$ or $q$'') is
true exactly when at least one of $p, q$ is true (inclusive or).
\end{definition}

\begin{definition}[Exclusive or]
The \emph{exclusive or} $p \oplus q$ is true exactly when $p$ and $q$ have
different truth values.
\end{definition}

\begin{definition}[Conditional]
The \emph{conditional} (material implication) $p \to q$ (``if $p$ then $q$'')
is false only when $p$ is true and $q$ is false; in all other cases it is true.
$p$ is the \emph{hypothesis} (antecedent), $q$ the \emph{conclusion}
(consequent).
\end{definition}

\begin{notebox}
A conditional with a false hypothesis is \emph{vacuously true}: ``if $1 = 2$
then I am the Pope'' is a true proposition. The conditional asserts nothing
about causality --- only that the case ($p$ true, $q$ false) does not occur.
\end{notebox}

\begin{definition}[Biconditional]
The \emph{biconditional} $p \leftrightarrow q$ (``$p$ if and only if $q$'') is
true exactly when $p$ and $q$ have the same truth value.
\end{definition}

\begin{definition}[NAND and NOR]
$p \uparrow q \equiv \neg(p \wedge q)$ (\emph{NAND}, Sheffer stroke) and
$p \downarrow q \equiv \neg(p \vee q)$ (\emph{NOR}, Peirce arrow).
\end{definition}

The complete behaviour of the binary connectives:

\begin{center}
\begin{tabular}{cc|cccccc}
\toprule
$p$ & $q$ & $p\wedge q$ & $p\vee q$ & $p\oplus q$ & $p\to q$ &
$p\leftrightarrow q$ & $p\uparrow q$\\
\midrule
\TT & \TT & \TT & \TT & \FF & \TT & \TT & \FF\\
\TT & \FF & \FF & \TT & \TT & \FF & \FF & \TT\\
\FF & \TT & \FF & \TT & \TT & \TT & \FF & \TT\\
\FF & \FF & \FF & \FF & \FF & \TT & \TT & \TT\\
\bottomrule
\end{tabular}
\end{center}

\section{Truth tables}

To build the truth table of a compound proposition with $n$ atomic letters,
list all $2^n$ valuations (conventionally from $\TT\cdots\TT$ down to
$\FF\cdots\FF$), then evaluate subformulas from the innermost outward.

\begin{example}
Truth table of $(p \to q) \wedge (\neg p \to \neg q)$:
\begin{center}
\begin{tabular}{cc|ccc|c}
\toprule
$p$ & $q$ & $p\to q$ & $\neg p$ & $\neg p \to \neg q$ & whole formula\\
\midrule
\TT & \TT & \TT & \FF & \TT & \TT\\
\TT & \FF & \FF & \FF & \TT & \FF\\
\FF & \TT & \TT & \TT & \FF & \FF\\
\FF & \FF & \TT & \TT & \TT & \TT\\
\bottomrule
\end{tabular}
\end{center}
The formula is true exactly when $p$ and $q$ agree --- it is equivalent to
$p \leftrightarrow q$.
\end{example}

\subsection{Precedence}

To limit parentheses, the standard precedence (strongest first) is:
$\neg$, then $\wedge$, then $\vee$, then $\to$, then $\leftrightarrow$.
Thus $\neg p \vee q \wedge r \to s$ reads
$\bigl((\neg p) \vee (q \wedge r)\bigr) \to s$. The conditional associates to
the right: $p \to q \to r$ means $p \to (q \to r)$. When in doubt,
parenthesize.

\section{Converse, inverse, contrapositive}

\begin{definition}
For the conditional $p \to q$:
the \emph{converse} is $q \to p$;
the \emph{inverse} is $\neg p \to \neg q$;
the \emph{contrapositive} is $\neg q \to \neg p$.
\end{definition}

\begin{theorem}
A conditional is logically equivalent to its contrapositive:
$p \to q \;\equiv\; \neg q \to \neg p$. It is \emph{not} in general equivalent
to its converse or to its inverse (which are equivalent to each other).
\end{theorem}

Two classical fallacies arise from confusing these forms:
\emph{affirming the consequent} (from $p \to q$ and $q$, wrongly concluding
$p$) and \emph{denying the antecedent} (from $p \to q$ and $\neg p$, wrongly
concluding $\neg q$).

\section{Tautology, contradiction, contingency}

\begin{definition}
A compound proposition is a \emph{tautology} if it is true under every
valuation; a \emph{contradiction} if it is false under every valuation; a
\emph{contingency} if it is neither.
\end{definition}

\begin{example}
$p \vee \neg p$ is a tautology (law of excluded middle); $p \wedge \neg p$ is a
contradiction; $p \to q$ is a contingency.
\end{example}

\begin{definition}[Satisfiability]
A formula is \emph{satisfiable} if at least one valuation makes it true. The
\emph{SAT problem} asks whether a given formula is satisfiable; it is the
canonical NP-complete problem, and a formula with $n$ variables has $2^n$
candidate valuations --- brute force quickly becomes infeasible.
\end{definition}

\section{Logical equivalence}

\begin{definition}
Two formulas $\varphi$ and $\psi$ are \emph{logically equivalent}, written
$\varphi \equiv \psi$, if they have the same truth value under every valuation
--- equivalently, if $\varphi \leftrightarrow \psi$ is a tautology.
\end{definition}

\begin{theorem}[Major equivalence laws]\label{thm:laws}
For all formulas $p, q, r$ (with $\TT$, $\FF$ denoting a tautology and a
contradiction):
\begin{center}
\begin{tabular}{ll}
\toprule
Identity & $p \wedge \TT \equiv p$, \quad $p \vee \FF \equiv p$\\
Domination & $p \vee \TT \equiv \TT$, \quad $p \wedge \FF \equiv \FF$\\
Idempotence & $p \vee p \equiv p$, \quad $p \wedge p \equiv p$\\
Double negation & $\neg(\neg p) \equiv p$\\
Commutativity & $p \vee q \equiv q \vee p$, \quad $p \wedge q \equiv q \wedge p$\\
Associativity & $(p \vee q) \vee r \equiv p \vee (q \vee r)$, and dually for $\wedge$\\
Distributivity & $p \vee (q \wedge r) \equiv (p \vee q) \wedge (p \vee r)$, and dually\\
De Morgan & $\neg(p \wedge q) \equiv \neg p \vee \neg q$, \quad
             $\neg(p \vee q) \equiv \neg p \wedge \neg q$\\
Absorption & $p \vee (p \wedge q) \equiv p$, \quad $p \wedge (p \vee q) \equiv p$\\
Negation & $p \vee \neg p \equiv \TT$, \quad $p \wedge \neg p \equiv \FF$\\
\bottomrule
\end{tabular}
\end{center}
\end{theorem}

\begin{theorem}[Implication and biconditional identities]\label{thm:impl}
\[
p \to q \equiv \neg p \vee q, \qquad
\neg(p \to q) \equiv p \wedge \neg q, \qquad
p \leftrightarrow q \equiv (p \to q) \wedge (q \to p).
\]
\end{theorem}

Each law is proved by a truth table; once proved, laws may be chained to
simplify formulas algebraically, replacing equals by equals.

\section{Logical consequence}

\begin{definition}[Entailment]
A set of formulas $\Gamma$ \emph{entails} $\varphi$, written
$\Gamma \models \varphi$, if every valuation making all formulas of $\Gamma$
true also makes $\varphi$ true. An \emph{argument} with premises $\Gamma$ and
conclusion $\varphi$ is \emph{valid} exactly when $\Gamma \models \varphi$.
\end{definition}

Validity can always be checked by a truth table, but with many atoms this is
exponential; Chapter~5 develops \emph{rules of inference} as a scalable
alternative. The most important rules --- modus ponens (from $p$ and
$p \to q$, infer $q$), modus tollens (from $\neg q$ and $p \to q$, infer
$\neg p$), hypothetical syllogism, disjunctive syllogism, addition,
simplification, conjunction and resolution --- are all entailments and are
treated in detail there.

\section{Normal forms}

\begin{definition}
A \emph{literal} is an atom or a negated atom. A formula is in
\emph{negation normal form} (NNF) if negations occur only on atoms.
It is in \emph{disjunctive normal form} (DNF) if it is a disjunction of
conjunctions of literals, and in \emph{conjunctive normal form} (CNF) if it is
a conjunction of disjunctions of literals.
\end{definition}

\begin{theorem}
Every propositional formula is equivalent to one in CNF and to one in DNF.
\end{theorem}

\begin{example}[CNF conversion]\label{ex:cnf}
Convert $\neg(p \to q) \vee r$ to CNF:
\[
\neg(p \to q) \vee r
\;\equiv\; (p \wedge \neg q) \vee r
\;\equiv\; (p \vee r) \wedge (\neg q \vee r),
\]
using Theorem~\ref{thm:impl} and then distributivity. The general algorithm:
eliminate $\leftrightarrow$ and $\to$, push negations inward by De Morgan
(reaching NNF), then distribute $\vee$ over $\wedge$.
\end{example}

A DNF can also be read directly off a truth table: each row where the formula
is true contributes one conjunction of literals describing that row.

\section{Functional completeness}

\begin{definition}
A set of connectives is \emph{functionally complete} if every Boolean function
$\B^n \to \B$ can be expressed using only connectives from the set.
\end{definition}

\begin{theorem}
$\{\neg, \wedge, \vee\}$ is functionally complete (via DNF); so are
$\{\neg, \wedge\}$ and $\{\neg, \vee\}$ (by De Morgan); and remarkably each of
the singletons $\{\uparrow\}$ and $\{\downarrow\}$ is functionally complete.
\end{theorem}

For NAND: $\neg p \equiv p \uparrow p$ and
$p \wedge q \equiv (p \uparrow q) \uparrow (p \uparrow q)$, which suffices.
This is why entire processors can be built from NAND gates alone.

\section{Logic gates and circuits}

Each connective corresponds to a hardware \emph{gate} (AND, OR, NOT, XOR,
NAND, NOR). Compound formulas correspond to combinational circuits. The
\emph{half-adder}, which adds two bits, illustrates the correspondence: the sum
bit is $s = p \oplus q$ and the carry bit is $c = p \wedge q$. Cascading
half-adders yields full adders and ultimately arithmetic units --- arithmetic
reduced to logic.

\section{Translating English to logic}

Useful patterns, with $p$ ``it rains'' and $q$ ``I take an umbrella'':

\begin{itemize}
\item ``If it rains, I take an umbrella'': $p \to q$.
\item ``I take an umbrella only if it rains'': $q \to p$.
\item ``I take an umbrella unless it rains'': $\neg p \to q$.
\item ``Taking an umbrella is necessary for rain'': $p \to q$
  ($q$ is the necessary condition).
\item ``Taking an umbrella is sufficient for rain'': $q \to p$.
\item ``I take an umbrella if and only if it rains'': $q \leftrightarrow p$.
\end{itemize}

\begin{notebox}
``Only if'' introduces the \emph{conclusion}, not the hypothesis; ``unless''
means ``if not''. These two patterns cause most translation errors.
\end{notebox}

\section{Summary}

Propositions take values in $\B = \{\TT,\FF\}$; connectives
$\neg, \wedge, \vee, \oplus, \to, \leftrightarrow, \uparrow, \downarrow$
combine them. Truth tables decide every semantic question in principle;
equivalence laws (Theorem~\ref{thm:laws}) decide them algebraically. Every
formula has a CNF and a DNF; $\{\uparrow\}$ alone suffices to express
everything; and logic realised in silicon becomes digital circuitry.

\section*{Exercises}
\addcontentsline{toc}{section}{Exercises}

\begin{exercise}{\easy}
Which of the following are propositions?
(a) ``$7$ is prime.''\quad (b) ``$x^2 \geq 0$.''\quad
(c) ``Read this chapter twice.''\quad (d) ``$1+1=3$.''\quad
(e) ``This statement is false.''
\end{exercise}

\begin{exercise}{\easy}
Let $p$: ``the file compiles'' and $q$: ``the tests pass''. Express in symbols:
(a) the file compiles but the tests fail;
(b) if the file does not compile then the tests fail;
(c) the tests pass only if the file compiles;
(d) the file compiles if and only if the tests pass.
\end{exercise}

\begin{exercise}{\easy}
Build the truth table of $(p \vee q) \wedge \neg(p \wedge q)$ and identify
which standard connective it represents.
\end{exercise}

\begin{exercise}{\medium}
State the converse, inverse and contrapositive of: ``If $n$ is divisible by
$6$, then $n$ is divisible by $3$.'' Which of the four statements are true for
all integers $n$?
\end{exercise}

\begin{exercise}{\medium}
Classify each formula as tautology, contradiction or contingency:
(a) $(p \to q) \vee (q \to p)$;\quad
(b) $(p \wedge \neg p) \to q$;\quad
(c) $(p \to q) \to p$.
\end{exercise}

\begin{exercise}{\medium}
Using the equivalence laws (not truth tables), simplify
$\neg(p \vee \neg q) \vee (\neg p \wedge \neg q)$ as far as possible, citing
each law used.
\end{exercise}

\begin{exercise}{\medium}
Prove using equivalence laws that
$(p \to r) \wedge (q \to r) \;\equiv\; (p \vee q) \to r$.
\end{exercise}

\begin{exercise}{\medium}
Convert $(p \wedge q) \to r$ to CNF, then to DNF, showing each step.
\end{exercise}

\begin{exercise}{\hard}
Show that $\{\to, \FF\}$ is functionally complete, where $\FF$ is the constant
false. \emph{Hint: first express $\neg p$, then $\vee$, then $\wedge$.}
\end{exercise}

\begin{exercise}{\hard}
Express $p \vee q$ and $p \to q$ using only the NOR connective $\downarrow$.
\end{exercise}

\begin{exercise}{\hard}
A formula has the truth table that is true exactly on the rows
$(p,q,r) = (\TT,\TT,\FF)$ and $(\FF,\TT,\TT)$. Write its DNF, and give an
equivalent CNF.
\end{exercise}

\begin{exercise}{\hard}
Three suspects make statements. Alice: ``Bob is lying.'' Bob: ``Carol is
lying.'' Carol: ``Alice and Bob are both lying.'' Using propositional
variables $a, b, c$ for ``Alice (resp.\ Bob, Carol) tells the truth'',
formalise the situation and determine who is telling the truth.
\end{exercise}

\section*{Solutions to the Exercises}
\addcontentsline{toc}{section}{Solutions to the Exercises}

\begin{solution}{2.1}
(a) Yes (true). (b) No as stated --- $x$ is a free variable, so it is a
predicate; (it becomes a proposition once quantified, e.g.\
$\forall x \in \R,\ x^2 \geq 0$). (c) No --- imperative. (d) Yes (false).
(e) No --- it can receive no consistent truth value.
\end{solution}

\begin{solution}{2.2}
(a) $p \wedge \neg q$. (b) $\neg p \to \neg q$. (c) $q \to p$.
(d) $p \leftrightarrow q$.
\end{solution}

\begin{solution}{2.3}
\begin{center}
\begin{tabular}{cc|ccc}
\toprule
$p$&$q$&$p\vee q$&$\neg(p\wedge q)$&formula\\
\midrule
\TT&\TT&\TT&\FF&\FF\\
\TT&\FF&\TT&\TT&\TT\\
\FF&\TT&\TT&\TT&\TT\\
\FF&\FF&\FF&\TT&\FF\\
\bottomrule
\end{tabular}
\end{center}
True exactly when $p$ and $q$ differ: this is the exclusive or $p \oplus q$.
\end{solution}

\begin{solution}{2.4}
Converse: ``if $n$ is divisible by $3$ then $n$ is divisible by $6$'' ---
false ($n = 9$). Inverse: ``if $n$ is not divisible by $6$ then $n$ is not
divisible by $3$'' --- false ($n = 9$ again). Contrapositive: ``if $n$ is not
divisible by $3$ then $n$ is not divisible by $6$'' --- true, being equivalent
to the original true statement.
\end{solution}

\begin{solution}{2.5}
(a) Tautology: if $p \to q$ fails then $p$ is true and $q$ false, making
$q \to p$ true. (b) Tautology: the hypothesis $p \wedge \neg p$ is always
false, so the conditional is vacuously true. (c) Contingency: true when
$p = \TT$; false when $p = \FF$ (then $p \to q$ is true but the conclusion
$p$ is false).
\end{solution}

\begin{solution}{2.6}
\begin{align*}
\neg(p \vee \neg q) \vee (\neg p \wedge \neg q)
&\equiv (\neg p \wedge q) \vee (\neg p \wedge \neg q) && \text{De Morgan, double negation}\\
&\equiv \neg p \wedge (q \vee \neg q) && \text{distributivity}\\
&\equiv \neg p \wedge \TT && \text{negation law}\\
&\equiv \neg p && \text{identity law.}
\end{align*}
\end{solution}

\begin{solution}{2.7}
\begin{align*}
(p \to r) \wedge (q \to r)
&\equiv (\neg p \vee r) \wedge (\neg q \vee r) && \text{implication identity}\\
&\equiv (\neg p \wedge \neg q) \vee r && \text{distributivity}\\
&\equiv \neg(p \vee q) \vee r && \text{De Morgan}\\
&\equiv (p \vee q) \to r && \text{implication identity.}
\end{align*}
\end{solution}

\begin{solution}{2.8}
$(p \wedge q) \to r \equiv \neg(p \wedge q) \vee r
\equiv \neg p \vee \neg q \vee r$. This single clause is simultaneously a CNF
(one disjunctive clause) and a DNF (disjunction of three one-literal
conjunctions).
\end{solution}

\begin{solution}{2.9}
$\neg p \equiv p \to \FF$. Then
$p \vee q \equiv \neg p \to q \equiv (p \to \FF) \to q$, and
$p \wedge q \equiv \neg(\neg p \vee \neg q)
\equiv \bigl(((p \to \FF) \to (q \to \FF))\bigr) \to \FF$ after expanding
$\neg p \vee \neg q \equiv p \to \neg q$, i.e.\
$p \wedge q \equiv (p \to (q \to \FF)) \to \FF$. Since
$\{\neg, \wedge, \vee\}$ is functionally complete, so is $\{\to, \FF\}$.
\end{solution}

\begin{solution}{2.10}
$\neg x \equiv x \downarrow x$. Then
$p \vee q \equiv \neg(p \downarrow q)
\equiv (p \downarrow q) \downarrow (p \downarrow q)$, and
$p \to q \equiv \neg p \vee q$, so with $A = p \downarrow p$ ($= \neg p$):
$p \to q \equiv (A \downarrow q) \downarrow (A \downarrow q)$.
\end{solution}

\begin{solution}{2.11}
DNF: each true row becomes a conjunction:
$(p \wedge q \wedge \neg r) \vee (\neg p \wedge q \wedge r)$.
For a CNF, note both rows have $q = \TT$, and $p \leftrightarrow \neg r$ holds
on them ($p$ and $r$ differ); conversely $q \wedge (p \oplus r)$ is true
exactly on these two rows. Hence a CNF is
$q \wedge (p \vee r) \wedge (\neg p \vee \neg r)$.
\end{solution}

\begin{solution}{2.12}
The statements give $a \leftrightarrow \neg b$, $b \leftrightarrow \neg c$,
$c \leftrightarrow (\neg a \wedge \neg b)$. From the first two: $a$ and $b$
differ, $b$ and $c$ differ, so $a = c$ and $a \neq b$. Case $a = \TT$: then
$c = \TT$, but Carol's statement requires $\neg a \wedge \neg b$, i.e.\
$a = \FF$ --- contradiction. Case $a = \FF$: then $b = \TT$, $c = \FF$. Check
Carol: she lies, and indeed $\neg a \wedge \neg b = \FF$ --- consistent.
So Bob alone tells the truth.
\end{solution}
